\documentclass[12pt,a4paper]{report}

\usepackage[T1]{fontenc}
\usepackage[utf8]{inputenc}
\usepackage[french]{babel}
\usepackage{geometry}
\usepackage{graphicx}
\usepackage{xcolor}
\usepackage{booktabs}
\usepackage{longtable}
\usepackage{array}
\usepackage{tabularx}
\usepackage{enumitem}
\usepackage{hyperref}
\usepackage{listings}
\usepackage{caption}
\usepackage{float}
\usepackage{tikz}
\usetikzlibrary{arrows.meta,positioning}

\geometry{left=2.5cm,right=2.5cm,top=2.5cm,bottom=2.5cm}
\setlength{\parindent}{0pt}
\setlength{\parskip}{0.5em}

\definecolor{codeblue}{RGB}{0,92,170}
\definecolor{codegray}{RGB}{95,95,95}
\definecolor{codegreen}{RGB}{0,128,0}
\definecolor{backgray}{RGB}{247,247,247}
\definecolor{titlegreen}{RGB}{35,92,71}

\lstdefinestyle{dartstyle}{
  language=C,
  basicstyle=\ttfamily\small,
  keywordstyle=\color{codeblue}\bfseries,
  commentstyle=\color{codegray}\itshape,
  stringstyle=\color{codegreen},
  showstringspaces=false,
  breaklines=true,
  frame=single,
  rulecolor=\color{codegray},
  backgroundcolor=\color{backgray},
  numbers=left,
  numberstyle=\tiny\color{codegray},
  numbersep=8pt,
  captionpos=b,
  tabsize=2
}
\lstset{style=dartstyle}
\hypersetup{
  colorlinks=true,
  linkcolor=titlegreen,
  urlcolor=codeblue,
  pdftitle={Rapport - Gestion des Produits}
}

\renewcommand{\chaptername}{Chapitre}
\renewcommand{\contentsname}{Table des matières}

\begin{document}

% ===================== Page de garde =====================
\begin{titlepage}
\centering

\vspace*{0.5cm}
\textbf{\Large [NOM DE L'ÉTABLISSEMENT]}\\[0.2cm]
\textbf{[Filière / Département]}\\[1.2cm]

\rule{\textwidth}{0.8pt}\\[0.7cm]
{\fontsize{28}{34}\selectfont
\textbf{Rapport de Projet}\\[0.3cm]
}
{\Large Application Mobile de Gestion d'un Coffee-Restaurant}\\[0.7cm]
\rule{\textwidth}{0.8pt}

\vspace{1.5cm}

\begin{tikzpicture}
\node[draw=titlegreen, very thick, rounded corners=8pt,
      minimum width=11cm, minimum height=2.2cm, align=center,
      fill=titlegreen!5]
{\Huge\bfseries Gestion des Produits};
\end{tikzpicture}

\vfill

\begin{minipage}[t]{0.48\textwidth}
\textbf{Réalisé par :}\\
[Nom et prénom]\\
[Classe / Groupe]
\end{minipage}
\hfill
\begin{minipage}[t]{0.48\textwidth}
\textbf{Encadré par :}\\
[Nom de l'encadrant]
\end{minipage}

\vfill
\textbf{Année universitaire : 2025--2026}
\end{titlepage}

% ===================== Remerciements =====================
\chapter*{Remerciements}
\addcontentsline{toc}{chapter}{Remerciements}

Je tiens à remercier toute personne ayant contribué à la réalisation de ce projet.
Je remercie particulièrement [Nom de l'encadrant] pour son encadrement, ses conseils
et son accompagnement tout au long du travail.

Je remercie également les enseignants, les responsables de l'établissement ainsi que
mes camarades pour les échanges techniques qui ont permis d'améliorer la qualité de
ce travail.

Enfin, je remercie ma famille pour son soutien, sa patience et ses encouragements.

% ===================== Résumé =====================
\chapter*{Résumé}
\addcontentsline{toc}{chapter}{Résumé}

Ce rapport présente la conception et la réalisation du module de gestion des produits
au sein d'une application mobile Flutter pour un coffee-restaurant. Ce module permet
à l'administrateur d'ajouter, consulter, modifier et supprimer les produits du menu.
Il offre aussi la gestion des catégories, des suppléments, des prix en dinars tunisiens
(DT), des images et de la disponibilité.

Le projet repose sur une architecture organisée en trois niveaux : les modèles
de données, les fournisseurs utilisant \texttt{Provider} pour la gestion d'état,
et les dépôts de données basés sur Firebase Cloud Firestore. L'application met en œuvre
les opérations CRUD, la validation des formulaires, le filtrage, la recherche et
la synchronisation des données avec la base distante.

\textbf{Mots-clés :} Flutter, Dart, Firebase, Cloud Firestore, CRUD, Provider,
gestion de produits.

\tableofcontents
\listoffigures
\listoftables

% ===================== Introduction =====================
\chapter{Introduction générale}

\section{Contexte du projet}

La digitalisation des établissements de restauration permet d'améliorer la gestion
des commandes, du menu et de la relation avec les clients. Dans une application de
coffee-restaurant, le catalogue des produits constitue une donnée centrale : les
produits sont présentés aux clients, associés à des catégories et affectés à des
suppléments optionnels.

L'application réalisée avec Flutter utilise Firebase pour centraliser les données
et proposer une application multi-plateformes. Le responsable peut gérer les produits
à partir d'un espace administrateur, tandis que les clients peuvent consulter le menu.

\section{Objectifs du module}

Les objectifs principaux du module de gestion des produits sont :
\begin{itemize}[label=\textbullet]
  \item centraliser les informations du menu ;
  \item permettre l'ajout et la modification d'un produit ;
  \item supprimer un produit après confirmation ;
  \item consulter les caractéristiques détaillées d'un produit ;
  \item changer rapidement son état de disponibilité ;
  \item rechercher un produit ou filtrer les produits par catégorie ;
  \item assurer la persistance des données dans Firebase ;
  \item informer l'utilisateur en cas de succès ou d'erreur.
\end{itemize}

\section{Organisation du rapport}

Le premier chapitre présente le contexte et les objectifs du projet. Le deuxième
chapitre décrit les principes de gestion des produits et les opérations CRUD.
Le troisième chapitre présente l'architecture, la réalisation du code et les écrans.
Le quatrième chapitre traite de la liaison avec Firebase. Enfin, le dernier chapitre
présente les tests, les difficultés rencontrées et les perspectives d'évolution.

% ===================== Présentation générale =====================
\chapter{Gestion des produits}

\section{Description fonctionnelle}

La gestion des produits permet au responsable d'administrer le catalogue du
coffee-restaurant. Un produit possède les informations suivantes :
\begin{itemize}[label=\textbullet]
  \item un identifiant unique généré par Firestore ;
  \item un nom ;
  \item une catégorie ;
  \item une description ;
  \item un prix en dinars tunisiens ;
  \item une image optionnelle ;
  \item un état de disponibilité ;
  \item une liste de suppléments.
\end{itemize}

Un supplément contient uniquement un nom et un prix. Il peut être sélectionné
à la commande pour personnaliser le produit.

\section{Cas d'utilisation}

\begin{figure}[H]
\centering
\begin{tikzpicture}[node distance=1.8cm and 2.2cm,
actor/.style={circle,draw,fill=titlegreen!15,minimum size=1.2cm},
usecase/.style={draw,rounded corners,fill=gray!8,minimum width=4cm,
               minimum height=1cm,align=center},
arrow/.style={-{Stealth[length=2.5mm]},thick}]
\node[actor] (admin) {Responsable};
\node[usecase, right=of admin,yshift=2.7cm] (lister) {Lister les produits};
\node[usecase, right=of admin] (ajouter) {Ajouter un produit};
\node[usecase, right=of admin,yshift=-2.7cm] (modifier) {Modifier un produit};
\node[usecase, right=of ajouter,xshift=1.2cm,yshift=1.35cm] (supprimer) {Supprimer un produit};
\node[usecase, right=of modifier,xshift=1.2cm,yshift=-1.35cm] (dispo) {Changer la disponibilité};
\draw[arrow] (admin)--(lister);
\draw[arrow] (admin)--(ajouter);
\draw[arrow] (admin)--(modifier);
\draw[arrow] (admin)--(supprimer);
\draw[arrow] (admin)--(dispo);
\end{tikzpicture}
\caption{Diagramme simplifié des fonctionnalités de gestion des produits}
\label{fig:cas_utilisation}
\end{figure}

\section{Les opérations CRUD}

\subsection{Création}

Depuis l'écran de gestion du menu, le responsable appuie sur le bouton
\texttt{+} puis renseigne le formulaire. Les champs obligatoires sont le nom,
la catégorie, la description et le prix. Le prix doit être strictement positif.
Une image peut être choisie dans la galerie.

La création produit un objet \texttt{Product} puis appelle
\texttt{MenuProvider.addProduct}. Le dépôt crée un document dans la collection
\texttt{products} avec un identifiant généré automatiquement.

\subsection{Lecture}

Au démarrage de l'application, \texttt{loadMenu()} charge simultanément les catégories
et les produits depuis Firestore. Les données sont stockées dans
\texttt{MenuProvider}. La liste peut être filtrée par texte ou par catégorie sans
interroger à nouveau la base.

\subsection{Modification}

Le responsable ouvre le formulaire d'un produit existant. Les valeurs actuelles
sont préchargées dans les contrôleurs. Après validation, \texttt{updateProduct}
met à jour le document correspondant à \texttt{product.id}. Le champ
\texttt{createdAt} n'est pas modifié.

\subsection{Suppression}

Avant la suppression, une boîte de dialogue demande une confirmation. Après
validation, \texttt{deleteProduct(productId)} supprime le document Firestore.
Le produit est ensuite retiré de la liste locale afin de mettre immédiatement
l'interface à jour.

\subsection{Changement de disponibilité}

Le gestionnaire utilise un interrupteur sur chaque carte de produit. Lorsque
l'administrateur modifie l'état, \texttt{toggleAvailability} crée une copie du
produit avec \texttt{isAvailable} inversé. Une mise à jour locale optimiste est
effectuée avant l'écriture distante.

\section{Tableau récapitulatif des CRUD}

\begin{table}[H]
\centering
\caption{Correspondance entre les actions et les opérations Firestore}
\begin{tabularx}{\textwidth}{>{\bfseries}p{2.5cm}p{3cm}p{4.1cm}X}
\toprule
Opération & Action utilisateur & Fonction applicative & Opération Firestore \\
\midrule
Create & Ajouter un produit & \texttt{addProduct()} & \texttt{doc().set()} \\
Read & Charger/filtrer & \texttt{fetchProducts()} & \texttt{collection().get()} \\
Update & Modifier & \texttt{updateProduct()} & \texttt{doc(id).update()} \\
Delete & Supprimer & \texttt{deleteProduct()} & \texttt{doc(id).delete()} \\
\bottomrule
\end{tabularx}
\end{table}

% ===================== Réalisation =====================
\chapter{Réalisation technique}

\section{Technologies utilisées}

\begin{table}[H]
\centering
\caption{Outils et technologies utilisés}
\begin{tabularx}{\textwidth}{>{\bfseries}p{3cm}X}
\toprule
Technologie & Rôle \\
\midrule
Flutter & Développement de l'interface multi-plateformes \\
Dart & Langage de programmation principal \\
Provider & Gestion et diffusion de l'état applicatif \\
Cloud Firestore & Stockage NoSQL des produits et catégories \\
Firebase Core & Initialisation de Firebase \\
Image Picker & Sélection d'une image depuis la galerie \\
Material Design & Construction des composants d'interface \\
\bottomrule
\end{tabularx}
\end{table}

\section{Architecture de l'application}

Le module suit une séparation des responsabilités :
\begin{itemize}[label=\textbullet]
  \item \texttt{models} : représentation des données ;
  \item \texttt{data} : accès à Firebase et interface du dépôt ;
  \item \texttt{providers} : état applicatif et notifications de l'interface ;
  \item \texttt{screens} : écrans de gestion et formulaires ;
  \item \texttt{widgets} : composants visuels réutilisables.
\end{itemize}

\begin{figure}[H]
\centering
\begin{tikzpicture}[node distance=1.2cm,
box/.style={draw,rounded corners,minimum width=4.2cm,minimum height=1cm,
            align=center,fill=titlegreen!8},
arrow/.style={-{Stealth[length=2.5mm]},thick}]
\node[box] (ui) {Screens / Widgets};
\node[box, below=of ui] (provider) {MenuProvider};
\node[box, below=of provider] (contract) {MenuRepository};
\node[box, below=of contract] (firebase) {FirebaseMenuRepository};
\node[box, below=of firebase] (db) {Cloud Firestore};
\draw[arrow] (ui)--(provider);
\draw[arrow] (provider)--(contract);
\draw[arrow] (contract)--(firebase);
\draw[arrow] (firebase)--(db);
\end{tikzpicture}
\caption{Architecture en couches du module}
\label{fig:architecture}
\end{figure}

\section{Modèle Produit}

La classe \texttt{Product} centralise les champs persistés. Elle fournit
également \texttt{toMap()} pour la conversion vers Firestore et
\texttt{fromMap()} pour la reconstruction depuis un document.

\lstinputlisting[
  language=C,
  caption={Modèle \texttt{Product}},
  label={lst:product}
]{lib/features/admin/menu/models/product.dart}

\section{Contrat du dépôt}

Le dépôt expose une abstraction indépendante de la technologie Firebase.
Cette séparation facilite les tests et permet éventuellement de remplacer
FireStore par une autre source de données.

\lstinputlisting[
  language=C,
  caption={Interface \texttt{MenuRepository}},
  label={lst:repository}
]{lib/features/admin/menu/data/menu_repository.dart}

\section{Gestion de l'état avec Provider}

\texttt{MenuProvider} contient la liste des produits, les catégories, les filtres,
l'état de chargement et les éventuelles erreurs. Les méthodes CRUD modifient
les données locales puis notifiant les widgets avec \texttt{notifyListeners()}.

\lstinputlisting[
  language=C,
  caption={Méthodes CRUD du \texttt{MenuProvider}},
  firstline=63,
  lastline=112,
  label={lst:provider_crud}
]{lib/features/admin/menu/providers/menu_provider.dart}

\section{Écran de gestion des produits}

L'écran affiche une barre de recherche, un filtre de catégories et une liste
de cartes produits. Chaque carte contient les actions voir, modifier et supprimer,
ainsi qu'un interrupteur de disponibilité.

\lstinputlisting[
  language=C,
  caption={Actions de gestion depuis l'écran principal},
  firstline=25,
  lastline=68,
  label={lst:screen_actions}
]{lib/features/admin/menu/screens/menu_management_screen.dart}

\section{Carte d'affichage d'un produit}

Le widget \texttt{AdminProductCard} construit l'affichage d'un produit dans la
liste administrateur. Il affiche l'image, le nom, la catégorie, le prix et
l'état de disponibilité. Il contient également les actions de consultation,
de modification et de suppression.

\lstinputlisting[
  language=C,
  caption={Carte \texttt{AdminProductCard}},
  label={lst:admin_product_card}
]{lib/features/admin/menu/widgets/admin_product_card.dart}

\section{Fiche détaillée d'un produit}

La fonction \texttt{showProductDetailsSheet()} ouvre une fenêtre modale contenant
les informations complètes du produit sélectionné : image, nom, catégorie,
prix, description et disponibilité. Le même fichier fournit également la fiche
détaillée d'une catégorie et la liste de ses produits.

\lstinputlisting[
  language=C,
  caption={Fiche détaillée et feuille modale du produit},
  label={lst:menu_item_details_sheet}
]{lib/features/admin/menu/widgets/menu_item_details_sheet.dart}

\section{Gestion des suppléments}

Le widget \texttt{SupplementsEditor} permet d'ajouter ou de retirer les
suppléments associés à un produit. Chaque supplément contient un nom et un
prix. Les lignes saisies sont filtrées avant d'être transmises au formulaire
principal grâce à la fonction de rappel \texttt{onChanged}.

\lstinputlisting[
  language=C,
  caption={Éditeur de suppléments du produit},
  label={lst:supplements_editor}
]{lib/features/admin/menu/widgets/supplements_editor.dart}

\section{Formulaire d'ajout et de modification}

Le formulaire utilise une \texttt{GlobalKey<FormState>} pour valider les champs
avant l'enregistrement. Le même écran est utilisé pour l'ajout et la modification :
la présence de \texttt{widget.product} détermine le mode.

\lstinputlisting[
  language=C,
  caption={Logique d'ajout et de modification},
  firstline=66,
  lastline=111,
  label={lst:form_save}
]{lib/features/admin/menu/screens/product_form_screen.dart}

\section{Gestion de l'image}

L'image est sélectionnée avec \texttt{ImagePicker}. Elle est redimensionnée,
compressée puis convertie en chaîne Base64 avec le préfixe
\texttt{data:image/jpeg;base64,}. Cette valeur est enregistrée dans le champ
\texttt{imageUrl}. Une limite est appliquée pour rester sous la taille maximale
recommandée d'un document Firestore.

\lstinputlisting[
  language=C,
  caption={Sélection et conversion de l'image},
  firstline=48,
  lastline=69,
  label={lst:image}
]{lib/features/admin/menu/screens/product_form_screen.dart}

\section{Affichage des produits côté client}

\texttt{ClientMenuScreen} présente les produits disponibles aux clients. Les
produits indisponibles sont exclus de la liste, puis le résultat peut être
filtré par catégorie et par recherche sur le nom ou la description.

\lstinputlisting[
  language=C,
  caption={Écran du menu côté client},
  label={lst:client_menu_screen}
]{lib/features/client/menu/screens/client_menu_screen.dart}

\section{Carte produit côté client}

Le widget \texttt{ClientProductTile} affiche l'image, le nom, la description et
le prix d'un produit dans la liste du menu client. Un premier appui ouvre la
fiche détaillée du produit, tandis que le bouton d'ajout permet de placer le
produit dans le panier.

\lstinputlisting[
  language=C,
  caption={Carte réutilisable \texttt{ClientProductTile}},
  label={lst:client_product_tile}
]{lib/features/client/menu/widgets/client_product_tile.dart}

\section{Fiche détaillée côté client}

\texttt{ProductDetailScreen} affiche les informations détaillées du produit
sélectionné. L'utilisateur peut choisir les suppléments proposés, modifier la
quantité et ajouter le produit au panier. Cette même fiche peut être utilisée
en mode sélection lors de l'ajout d'un produit à une commande.

\lstinputlisting[
  language=C,
  caption={Fiche détaillée du produit côté client},
  label={lst:client_product_detail}
]{lib/features/client/menu/screens/product_detail_screen.dart}

% ===================== Firebase =====================
\chapter{Liaison avec Firebase}

\section{Configuration du projet}

Firebase Cloud Firestore est initialisé au démarrage de l'application dans
\texttt{main()}. La méthode \texttt{Firebase.initializeApp()} reçoit les options
générées pour la plateforme courante par FlutterFire CLI.

\lstinputlisting[
  language=C,
  caption={Initialisation de Firebase dans \texttt{main()}},
  firstline=21,
  lastline=39,
  label={lst:firebase_init}
]{lib/main.dart}

\section{Structure des données}

Le module utilise principalement deux collections :
\begin{itemize}[label=\textbullet]
  \item \texttt{categories} : catégories du menu ;
  \item \texttt{products} : informations des produits.
\end{itemize}

\begin{table}[H]
\centering
\caption{Structure du document \texttt{products}}
\begin{tabularx}{\textwidth}{>{\bfseries}p{3cm}p{3cm}X}
\toprule
Champ & Type & Exemple \\
\midrule
name & String & Cappuccino \\
categoryId & String & ID du document catégorie \\
description & String & Café espresso avec lait mousseux \\
price & Number & 6.500 \\
imageUrl & String/null & URL, asset ou data URL \\
isAvailable & Boolean & true \\
supplements & Array & \texttt{[\{name: Noisette, price: 1\}]} \\
createdAt & Timestamp & Horodatage serveur \\
\bottomrule
\end{tabularx}
\end{table}

\section{Implémentation du dépôt Firebase}

\texttt{FirebaseMenuRepository} encapsule toutes les requêtes vers Firestore.
Il crée les références des collections \texttt{categories} et \texttt{products}
puis réalise les opérations CRUD.

\lstinputlisting[
  language=C,
  caption={Implémentation Firebase du dépôt menu},
  firstline=1,
  lastline=101,
  label={lst:firebase_repository}
]{lib/features/admin/menu/data/firebase_menu_repository.dart}

\section{Lecture des produits}

Les produits sont chargés avec \texttt{orderBy('createdAt')}. Cette approche permet
de conserver un ordre chronologique, tandis que \texttt{fromMap()} convertit chaque
document en objet Dart.

\section{Ajout d'un produit}

Lors de l'ajout, \texttt{\_products.doc()} crée une référence dotée d'un identifiant
généré par Firestore. L'opération \texttt{set()} écrit les données du produit et
ajoute \texttt{createdAt} avec \texttt{FieldValue.serverTimestamp()}.

\section{Modification d'un produit}

La modification utilise \texttt{update()} sur le document identifié par
\texttt{product.id}. Cette méthode ne remplace pas le champ \texttt{createdAt},
qui reste inchangé.

\section{Suppression d'un produit}

La suppression utilise \texttt{delete()} sur le document ciblé. Elle est précédée
dans l'interface par une confirmation afin d'éviter une suppression involontaire.

\section{Injection du dépôt}

Au démarrage, \texttt{FirebaseMenuRepository()} est injecté dans
\texttt{CoffeeShopApp}. \texttt{MenuProvider} utilise ensuite cette implémentation
pour charger les données et effectuer les opérations CRUD.

\begin{lstlisting}[language=C,caption={Partie du gestionnaire des produits dans \texttt{main.dart}}]
ChangeNotifierProvider<MenuProvider>(
  // Voici la partie de gestion des produits :
  // FirebaseMenuRepository fournit les données de MenuProvider.
  create: (_) => MenuProvider(menuRepository)..loadMenu(),
),
\end{lstlisting}

\lstinputlisting[
  language=C,
  caption={Injection du dépôt Firebase dans Provider},
  firstline=41,
  lastline=68,
  label={lst:provider_injection}
]{lib/main.dart}

% ===================== Tests / conclusion =====================
\chapter{Tests et conclusion}

\section{Scénarios de test}

\begin{table}[H]
\centering
\caption{Principaux scénarios testés}
\begin{tabularx}{\textwidth}{>{\bfseries}p{2.8cm}p{4.5cm}X}
\toprule
Scénario & Test effectué & Résultat attendu \\
\midrule
Ajout valide & Saisie de tous les champs obligatoires & Produit ajouté à Firestore \\
Champ manquant & Nom ou catégorie vide & Message de validation \\
Prix invalide & Prix nul ou texte & Enregistrement refusé \\
Modification & Changement du nom ou du prix & Document Firestore mis à jour \\
Suppression & Confirmation puis annulation & Document supprimé ou conservé \\
Disponibilité & Activation/désactivation & Champ modifié dans Firestore \\
Recherche & Saisie du nom ou de la description & Liste filtrée localement \\
Filtre catégorie & Sélection d'une catégorie & Produits associés seulement \\
Erreur réseau & Arrêt temporaire de Firebase & Message d'erreur affiché \\
\bottomrule
\end{tabularx}
\end{table}

\section{Résultats obtenus}

Le module permet de gérer l'ensemble du cycle de vie d'un produit. Les opérations
de lecture, création, modification et suppression communiquent avec Firebase.
La recherche et le filtre améliorent l'expérience de l'administrateur.

L'état \texttt{isAvailable} permet de masquer temporairement un produit aux clients
sans supprimer sa fiche. Les produits sont associés à une catégorie et peuvent
comprendre plusieurs suppléments.

\section{Difficultés rencontrées}

Plusieurs difficultés ont été rencontrées lors du développement :
\begin{itemize}[label=\textbullet]
  \item gestion de la conversion entre les objets Dart et les documents Firestore ;
  \item conservation des identifiants des documents ;
  \item validation des valeurs numériques ;
  \item gestion des erreurs réseau ;
  \item limitation de la taille des images enregistrées dans Firestore ;
  \item synchronisation de l'état local après chaque opération.
\end{itemize}

\section{Conclusion générale}

Le module de gestion des produits réalisé avec Flutter et Firebase répond aux
besoins essentiels d'administration du menu. Il fournit des interfaces claires,
des validations, une gestion centralisée de l'état et une persistance distante.

L'architecture proposée sépare l'interface, la logique métier et l'accès aux données.
Cette organisation rend le code plus lisible, maintenable et évolutif.

\section{Perspectives}

Plusieurs améliorations peuvent être envisagées :
\begin{itemize}[label=\textbullet]
  \item utiliser Firebase Storage pour les images ;
  \item ajouter les règles de sécurité Firestore ;
  \item afficher les données en temps réel avec des \texttt{StreamBuilder} ;
  \item ajouter des tests unitaires et des tests d'intégration ;
  \item ajouter une pagination pour les grands catalogues ;
  \item gérer des rôles plus détaillés pour les administrateurs ;
  \item ajouter une gestion des promotions et des variantes.
\end{itemize}

% ===================== Annexes =====================
\appendix
\chapter{Extrait du modèle complémentaire}

\lstinputlisting[
  language=C,
  caption={Modèle \texttt{Supplement}},
  label={lst:supplement}
]{lib/features/admin/menu/models/supplement.dart}

\chapter{Extrait du modèle catégorie}

\lstinputlisting[
  language=C,
  caption={Modèle \texttt{Category}},
  label={lst:category}
]{lib/features/admin/menu/models/category.dart}

\begin{thebibliography}{9}
\bibitem{flutter}
Flutter Documentation.
\newblock \url{https://docs.flutter.dev/}

\bibitem{firebase}
Firebase Documentation.
\newblock \url{https://firebase.google.com/docs}

\bibitem{firestore}
Cloud Firestore Documentation.
\newblock \url{https://firebase.google.com/docs/firestore}

\bibitem{provider}
Provider Package.
\newblock \url{https://pub.dev/packages/provider}
\end{thebibliography}

\end{document}
